\documentclass[paper=a4paper,fontsize=10pt,jafontsize=9pt,titlepage]{jlreq}
\usepackage{pxrubrica} %ルビ、傍点
\usepackage{longtable}
\usepackage[polutonikogreek,english,japanese]{babel}
\usepackage{latexsym}
\usepackage{color}
\usepackage{url}

%----- header -------
\usepackage{fancyhdr}
\pagestyle{fancy}
\lhead{{\itshape Summa Theologiae} I, q.79}
%--------------------


\bibliographystyle{jplain}

\title{{\bfseries Prima Pars}\\{\Huge Summae Theologiae}\\Sancti Thomae
Aquinatis\\{\sffamily Quaestio Septuagesimanona}\\{\bfseries De Potentiis Intellectivis}}
\author{Japanese translation\\by Yoshinori {\scshape Ueeda}}
\date{Last modified \today}

%%%% コピペ用
%\rhead{a.~}
%\begin{center}
% {\Large {\bfseries }}\\
% {\large }\\
% {\footnotesize }\\
% {\Large \\}
%\end{center}
%
%\begin{longtable}{p{21em}p{21em}}
%
%&
%
%\\
%\end{longtable}
%\newpage

\begin{document}

\maketitle

\begin{center}
{\Large 第七十九問\\
知性的能力について}
\end{center}


\begin{longtable}{p{21em}p{21em}}
 Deinde quaeritur de potentiis intellectivis. Circa quod quaeruntur tredecim.

 &


 次に、知性的能力について問われる。これを巡って、十三のことが問われる。
 
\\[1em]

 \begin{enumerate}
  \item utrum intellectus sit potentia animae, vel eius essentia.
  \item si est potentia, utrum sit potentia passiva.
  \item si est potentia passiva, utrum sit ponere aliquem intellectum agentem.
  \item utrum sit aliquid animae.
  \item utrum intellectus agens sit unus omnium.
  \item utrum memoria sit in intellectu.
  \item utrum sit alia potentia ab intellectu.
  \item utrum ratio sit alia potentia ab intellectu.
  \item utrum ratio superior et inferior sint diversae potentiae.
  \item utrum intelligentia sit alia potentia praeter intellectum.
  \item utrum intellectus speculativus et practicus sint diversae potentiae.
  \item utrum synderesis sit aliqua potentia intellectivae partis.
  \item utrum conscientia sit aliqua potentia intellectivae partis.
 \end{enumerate}

&
 
\begin{enumerate}
 \item 知性は魂の能力か、それとも本質か。
 \item もし能力ならば、それは受動的な能力か。
 \item もし受動的な能力なら、何か能動する知性を措定することができるか。
 \item それは魂に属する何かか。
 \item 能動知性は万人に一つか。
 \item 記憶が知性の中にあるか。
 \item 記憶は知性とは別の能力か。
 \item 理性は知性と別の能力か。
 \item 上位の理性と解の理性は、異なる能力に属するか。
 \item 知解は知性と別の能力か。
 \item 観照的知性と実践的知性は別の能力か。
 \item 良知は知性的部分の何らかの能力か。
 \item 良心は知性的部分の何らかの能力か。
\end{enumerate}
\end{longtable}

\newpage

\end{document}

QUAESTIO 79
PROOEMIUM
[31715] Iª q. 79 pr.



ARTICULUS 1
[31716] Iª q. 79 a. 1 arg. 1
Ad primum sic proceditur. Videtur quod intellectus non sit aliqua potentia animae, sed sit ipsa eius essentia. Intellectus idem enim videtur esse quod mens. Sed mens non est potentia animae sed essentia, dicit enim Augustinus, IX de Trin., mens et spiritus non relative dicuntur, sed essentiam demonstrant. Ergo intellectus est ipsa essentia animae.

[31717] Iª q. 79 a. 1 arg. 2
Praeterea, diversa genera potentiarum animae non uniuntur in aliqua potentia una, sed in sola essentia animae. Appetitivum autem et intellectivum sunt diversa genera potentiarum animae, ut dicitur in II de anima; conveniunt autem in mente, quia Augustinus, X de Trin., ponit intelligentiam et voluntatem in mente. Ergo mens et intellectus est ipsa essentia animae, et non aliqua eius potentia.

[31718] Iª q. 79 a. 1 arg. 3
Praeterea, secundum Gregorium, in homilia ascensionis, homo intelligit cum Angelis. Sed Angeli dicuntur mentes et intellectus. Ergo mens et intellectus hominis non est aliqua potentia animae, sed ipsa anima.

[31719] Iª q. 79 a. 1 arg. 4
Praeterea, ex hoc convenit alicui substantiae quod sit intellectiva, quia est immaterialis. Sed anima est immaterialis per suam essentiam. Ergo videtur quod anima per suam essentiam sit intellectiva.

[31720] Iª q. 79 a. 1 s. c.
Sed contra est quod philosophus ponit intellectivum potentiam animae, ut patet in II de anima.

[31721] Iª q. 79 a. 1 co.
Respondeo dicendum quod necesse est dicere, secundum praemissa, quod intellectus sit aliqua potentia animae, et non ipsa animae essentia. Tunc enim solum immediatum principium operationis est ipsa essentia rei operantis, quando ipsa operatio est eius esse, sicut enim potentia se habet ad operationem ut ad suum actum, ita se habet essentia ad esse. In solo Deo autem idem est intelligere quod suum esse. Unde in solo Deo intellectus est eius essentia, in aliis autem creaturis intellectualibus intellectus est quaedam potentia intelligentis.

[31722] Iª q. 79 a. 1 ad 1
Ad primum ergo dicendum quod sensus accipitur aliquando pro potentia, aliquando vero pro ipsa anima sensitiva, denominatur enim anima sensitiva nomine principalioris suae potentiae, quae est sensus. Et similiter anima intellectiva quandoque nominatur nomine intellectus, quasi a principaliori sua virtute; sicut dicitur in I de anima, quod intellectus est substantia quaedam. Et etiam hoc modo Augustinus dicit quod mens est spiritus, vel essentia.

[31723] Iª q. 79 a. 1 ad 2
Ad secundum dicendum quod appetitivum et intellectivum sunt diversa genera potentiarum animae, secundum diversas rationes obiectorum. Sed appetitivum partim convenit cum intellectivo, et partim cum sensitivo, quantum ad modum operandi per organum corporale, vel sine huiusmodi organo, nam appetitus sequitur apprehensionem. Et secundum hoc Augustinus ponit voluntatem in mente, et philosophus in ratione.

[31724] Iª q. 79 a. 1 ad 3
Ad tertium dicendum quod in Angelis non est alia vis nisi intellectiva, et voluntas, quae ad intellectum consequitur. Et propter hoc Angelus dicitur mens vel intellectus, quia tota virtus sua in hoc consistit. Anima autem habet multas alias vires, sicut sensitivas et nutritivas, et ideo non est simile.

[31725] Iª q. 79 a. 1 ad 4
Ad quartum dicendum quod ipsa immaterialitas substantiae intelligentis creatae non est eius intellectus; sed ex immaterialitate habet virtutem ad intelligendum. Unde non oportet quod intellectus sit substantia animae, sed eius virtus et potentia.


ARTICULUS 2
[31726] Iª q. 79 a. 2 arg. 1
Ad secundum sic proceditur. Videtur quod intellectus non sit potentia passiva. Patitur enim unumquodque secundum materiam; sed agit ratione formae. Sed virtus intellectiva consequitur immaterialitatem substantiae intelligentis. Ergo videtur quod intellectus non sit potentia passiva.

[31727] Iª q. 79 a. 2 arg. 2
Praeterea, potentia intellectiva est incorruptibilis, ut supra dictum est. Sed intellectus si est passivus, est corruptibilis, ut dicitur in III de anima. Ergo potentia intellectiva non est passiva.

[31728] Iª q. 79 a. 2 arg. 3
Praeterea, agens est nobilius patiente, ut dicit Augustinus XII super Gen. ad Litt., et Aristoteles in III de anima. Potentiae autem vegetativae partis omnes sunt activae, quae tamen sunt infimae inter potentias animae. Ergo multo magis potentiae intellectivae, quae sunt supremae, omnes sunt activae.

[31729] Iª q. 79 a. 2 s. c.
Sed contra est quod philosophus dicit, in III de anima, quod intelligere est pati quoddam.

[31730] Iª q. 79 a. 2 co.
Respondeo dicendum quod pati tripliciter dicitur. Uno modo, propriissime, scilicet quando aliquid removetur ab eo quod convenit sibi secundum naturam, aut secundum propriam inclinationem; sicut cum aqua frigiditatem amittit per calefactionem, et cum homo aegrotat aut tristatur. Secundo modo, minus proprie dicitur aliquis pati ex eo quod aliquid ab ipso abiicitur, sive sit ei conveniens, sive non conveniens. Et secundum hoc dicitur pati non solum qui aegrotat, sed etiam qui sanatur; non solum qui tristatur, sed etiam qui laetatur; vel quocumque modo aliquis alteretur vel moveatur. Tertio modo, dicitur aliquid pati communiter, ex hoc solo quod id quod est in potentia ad aliquid, recipit illud ad quod erat in potentia, absque hoc quod aliquid abiiciatur. Secundum quem modum, omne quod exit de potentia in actum, potest dici pati, etiam cum perficitur. Et sic intelligere nostrum est pati. Quod quidem hac ratione apparet. Intellectus enim, sicut supra dictum est, habet operationem circa ens in universali. Considerari ergo potest utrum intellectus sit in actu vel potentia, ex hoc quod consideratur quomodo intellectus se habeat ad ens universale. Invenitur enim aliquis intellectus qui ad ens universale se habet sicut actus totius entis, et talis est intellectus divinus, qui est Dei essentia, in qua originaliter et virtualiter totum ens praeexistit sicut in prima causa. Et ideo intellectus divinus non est in potentia, sed est actus purus. Nullus autem intellectus creatus potest se habere ut actus respectu totius entis universalis, quia sic oporteret quod esset ens infinitum. Unde omnis intellectus creatus, per hoc ipsum quod est, non est actus omnium intelligibilium, sed comparatur ad ipsa intelligibilia sicut potentia ad actum. Potentia autem dupliciter se habet ad actum. Est enim quaedam potentia quae semper est perfecta per actum; sicut diximus de materia corporum caelestium. Quaedam autem potentia est, quae non semper est in actu, sed de potentia procedit in actum; sicut invenitur in generabilibus et corruptibilibus. Intellectus igitur angelicus semper est in actu suorum intelligibilium, propter propinquitatem ad primum intellectum, qui est actus purus, ut supra dictum est. Intellectus autem humanus, qui est infimus in ordine intellectuum, et maxime remotus a perfectione divini intellectus, est in potentia respectu intelligibilium, et in principio est sicut tabula rasa in qua nihil est scriptum, ut philosophus dicit in III de anima. Quod manifeste apparet ex hoc, quod in principio sumus intelligentes solum in potentia, postmodum autem efficimur intelligentes in actu. Sic igitur patet quod intelligere nostrum est quoddam pati, secundum tertium modum passionis. Et per consequens intellectus est potentia passiva.

[31731] Iª q. 79 a. 2 ad 1
Ad primum ergo dicendum quod obiectio illa procedit de primo et secundo modo passionis, qui sunt proprii materiae primae. Tertius autem modus passionis est cuiuscumque in potentia existentis quod in actum reducitur.

[31732] Iª q. 79 a. 2 ad 2
Ad secundum dicendum quod intellectus passivus secundum quosdam dicitur appetitus sensitivus, in quo sunt animae passiones; qui etiam in I Ethic. dicitur rationalis per participationem, quia obedit rationi. Secundum alios autem intellectus passivus dicitur virtus cogitativa, quae nominatur ratio particularis. Et utroque modo passivum accipi potest secundum primos duos modos passionis, inquantum talis intellectus sic dictus, est actus alicuius organi corporalis. Sed intellectus qui est in potentia ad intelligibilia, quem Aristoteles ob hoc nominat intellectum possibilem, non est passivus nisi tertio modo, quia non est actus organi corporalis. Et ideo est incorruptibilis.

[31733] Iª q. 79 a. 2 ad 3
Ad tertium dicendum quod agens est nobilius patiente, si ad idem actio et passio referantur, non autem semper, si ad diversa. Intellectus autem est vis passiva respectu totius entis universalis. Vegetativum autem est activum respectu cuiusdam entis particularis, scilicet corporis coniuncti. Unde nihil prohibet huiusmodi passivum esse nobilius tali activo.


ARTICULUS 3
[31734] Iª q. 79 a. 3 arg. 1
Ad tertium sic proceditur. Videtur quod non sit ponere intellectum agentem. Sicut enim se habet sensus ad sensibilia, ita se habet intellectus noster ad intelligibilia. Sed quia sensus est in potentia ad sensibilia non ponitur sensus agens, sed sensus patiens tantum. Ergo, cum intellectus noster sit in potentia ad intelligibilia, videtur quod non debeat poni intellectus agens, sed possibilis tantum.

[31735] Iª q. 79 a. 3 arg. 2
Praeterea, si dicatur quod in sensu etiam est aliquod agens, sicut lumen, contra, lumen requiritur ad visum inquantum facit medium lucidum in actu, nam color ipse secundum se est motivus lucidi. Sed in operatione intellectus non ponitur aliquod medium quod necesse sit fieri in actu. Ergo non est necessarium ponere intellectum agentem.

[31736] Iª q. 79 a. 3 arg. 3
Praeterea, similitudo agentis recipitur in patiente secundum modum patientis. Sed intellectus possibilis est virtus immaterialis. Ergo immaterialitas eius sufficit ad hoc quod recipiantur in eo formae immaterialiter. Sed ex hoc ipso aliqua forma est intelligibilis in actu, quod est immaterialis. Ergo nulla necessitas est ponere intellectum agentem, ad hoc quod faciat species intelligibiles in actu.

[31737] Iª q. 79 a. 3 s. c.
Sed contra est quod philosophus dicit, in III de anima, quod sicut in omni natura ita et in anima est aliquid quo est omnia fieri, et aliquid quo est omnia facere. Est ergo ponere intellectum agentem.

[31738] Iª q. 79 a. 3 co.
Respondeo dicendum quod, secundum opinionem Platonis, nulla necessitas erat ponere intellectum agentem ad faciendum intelligibilia in actu; sed forte ad praebendum lumen intelligibile intelligenti, ut infra dicetur. Posuit enim Plato formas rerum naturalium sine materia subsistere, et per consequens eas intelligibiles esse, quia ex hoc est aliquid intelligibile actu, quod est immateriale. Et huiusmodi vocabat species, sive ideas, ex quarum participatione dicebat etiam materiam corporalem formari, ad hoc quod individua naturaliter constituerentur in propriis generibus et speciebus; et intellectus nostros, ad hoc quod de generibus et speciebus rerum scientiam haberent. Sed quia Aristoteles non posuit formas rerum naturalium subsistere sine materia; formae autem in materia existentes non sunt intelligibiles actu, sequebatur quod naturae seu formae rerum sensibilium, quas intelligimus, non essent intelligibiles actu. Nihil autem reducitur de potentia in actum, nisi per aliquod ens actu, sicut sensus fit in actu per sensibile in actu. Oportebat igitur ponere aliquam virtutem ex parte intellectus, quae faceret intelligibilia in actu, per abstractionem specierum a conditionibus materialibus. Et haec est necessitas ponendi intellectum agentem.

[31739] Iª q. 79 a. 3 ad 1
Ad primum ergo dicendum quod sensibilia inveniuntur actu extra animam, et ideo non oportuit ponere sensum agentem. Et sic patet quod in parte nutritiva omnes potentiae sunt activae; in parte autem sensitiva, omnes passivae; in parte vero intellectiva est aliquid activum, et aliquid passivum.

[31740] Iª q. 79 a. 3 ad 2
Ad secundum dicendum quod circa effectum luminis est duplex opinio. Quidam enim dicunt quod lumen requiritur ad visum, ut faciat colores actu visibiles. Et secundum hoc, similiter requiritur, et propter idem, intellectus agens ad intelligendum, propter quod lumen ad videndum. Secundum alios vero, lumen requiritur ad videndum, non propter colores, ut fiant actu visibiles; sed ut medium fiat actu lucidum, ut Commentator dicit in II de anima. Et secundum hoc, similitudo qua Aristoteles assimilat intellectum agentem lumini, attenditur quantum ad hoc, quod sicut hoc est necessarium ad videndum, ita illud ad intelligendum; sed non propter idem.

[31741] Iª q. 79 a. 3 ad 3
Ad tertium dicendum quod, supposito agente, bene contingit diversimode recipi eius similitudinem in diversis propter eorum dispositionem diversam. Sed si agens non praeexistit, nihil ad hoc faciet dispositio recipientis. Intelligibile autem in actu non est aliquid existens in rerum natura, quantum ad naturam rerum sensibilium, quae non subsistunt praeter materiam. Et ideo ad intelligendum non sufficeret immaterialitas intellectus possibilis, nisi adesset intellectus agens, qui faceret intelligibilia in actu per modum abstractionis.


ARTICULUS 4
[31742] Iª q. 79 a. 4 arg. 1
Ad quartum sic proceditur. Videtur quod intellectus agens non sit aliquid animae nostrae. Intellectus enim agentis effectus est illuminare ad intelligendum. Sed hoc fit per aliquid quod est altius anima; secundum illud Ioan. I, erat lux vera, quae illuminat omnem hominem venientem in hunc mundum. Ergo videtur quod intellectus agens non sit aliquid animae.

[31743] Iª q. 79 a. 4 arg. 2
Praeterea, philosophus, in III de anima, attribuit intellectui agenti quod non aliquando intelligit et aliquando non intelligit. Sed anima nostra non semper intelligit; sed aliquando intelligit et aliquando non intelligit. Ergo intellectus agens non est aliquid animae nostrae.

[31744] Iª q. 79 a. 4 arg. 3
Praeterea, agens et patiens sufficiunt ad agendum. Si igitur intellectus possibilis est aliquid animae nostrae, qui est virtus passiva, et similiter intellectus agens, qui est virtus activa; sequitur quod homo semper poterit intelligere cum voluerit, quod patet esse falsum. Non est ergo intellectus agens aliquid animae nostrae.

[31745] Iª q. 79 a. 4 arg. 4
Praeterea, philosophus dicit, in III de anima, quod intellectus agens est substantia actu ens. Nihil autem est respectu eiusdem in actu et in potentia. Si ergo intellectus possibilis, qui est in potentia ad omnia intelligibilia, est aliquid animae nostrae; videtur impossibile quod intellectus agens sit aliquid animae nostrae.

[31746] Iª q. 79 a. 4 arg. 5
Praeterea, si intellectus agens est aliquid animae nostrae, oportet quod sit aliqua potentia. Non est enim nec passio nec habitus, nam habitus et passiones non habent rationem agentis respectu passionum animae; sed magis passio est ipsa actio potentiae passivae, habitus autem est aliquid quod ex actibus consequitur. Omnis autem potentia fluit ab essentia animae. Sequeretur ergo quod intellectus agens ab essentia animae procederet. Et sic non inesset animae per participationem ab aliquo superiori intellectu, quod est inconveniens. Non ergo intellectus agens est aliquid animae nostrae.

[31747] Iª q. 79 a. 4 s. c.
Sed contra est quod philosophus dicit, III de anima quod necesse est in anima has esse differentias, scilicet intellectum possibilem, et agentem.

[31748] Iª q. 79 a. 4 co.
Respondeo dicendum quod intellectus agens de quo philosophus loquitur, est aliquid animae. Ad cuius evidentiam, considerandum est quod supra animam intellectivam humanam necesse est ponere aliquem superiorem intellectum, a quo anima virtutem intelligendi obtineat. Semper enim quod participat aliquid, et quod est mobile, et quod est imperfectum, praeexigit ante se aliquid quod est per essentiam suam tale, et quod est immobile et perfectum. Anima autem humana intellectiva dicitur per participationem intellectualis virtutis, cuius signum est, quod non tota est intellectiva, sed secundum aliquam sui partem. Pertingit etiam ad intelligentiam veritatis cum quodam discursu et motu, arguendo. Habet etiam imperfectam intelligentiam, tum quia non omnia intelligit; tum quia in his quae intelligit, de potentia procedit ad actum. Oportet ergo esse aliquem altiorem intellectum, quo anima iuvetur ad intelligendum. Posuerunt ergo quidam hunc intellectum secundum substantiam separatum, esse intellectum agentem, qui quasi illustrando phantasmata, facit ea intelligibilia actu. Sed, dato quod sit aliquis talis intellectus agens separatus, nihilominus tamen oportet ponere in ipsa anima humana aliquam virtutem ab illo intellectu superiori participatam, per quam anima humana facit intelligibilia in actu. Sicut et in aliis rebus naturalibus perfectis, praeter universales causas agentes, sunt propriae virtutes inditae singulis rebus perfectis, ab universalibus agentibus derivatae, non enim solus sol generat hominem, sed est in homine virtus generativa hominis; et similiter in aliis animalibus perfectis. Nihil autem est perfectius in inferioribus rebus anima humana. Unde oportet dicere quod in ipsa sit aliqua virtus derivata a superiori intellectu, per quam possit phantasmata illustrare. Et hoc experimento cognoscimus, dum percipimus nos abstrahere formas universales a conditionibus particularibus, quod est facere actu intelligibilia. Nulla autem actio convenit alicui rei, nisi per aliquod principium formaliter ei inhaerens; ut supra dictum est, cum de intellectu possibili ageretur. Ergo oportet virtutem quae est principium huius actionis, esse aliquid in anima. Et ideo Aristoteles comparavit intellectum agentem lumini, quod est aliquid receptum in aere. Plato autem intellectum separatum imprimentem in animas nostras, comparavit soli; ut Themistius dicit in commentario tertii de anima. Sed intellectus separatus, secundum nostrae fidei documenta, est ipse Deus, qui est creator animae, et in quo solo beatificatur, ut infra patebit. Unde ab ipso anima humana lumen intellectuale participat, secundum illud Psalmi IV, signatum est super nos lumen vultus tui, domine.

[31749] Iª q. 79 a. 4 ad 1
Ad primum ergo dicendum quod illa lux vera illuminat sicut causa universalis, a qua anima humana participat quandam particularem virtutem, ut dictum est.

[31750] Iª q. 79 a. 4 ad 2
Ad secundum dicendum quod philosophus illa verba non dicit de intellectu agente, sed de intellectu in actu. Unde supra de ipso praemiserat, idem autem est secundum actum scientia rei. Vel, si intelligatur de intellectu agente, hoc dicitur quia non est ex parte intellectus agentis hoc quod quandoque intelligimus et quandoque non intelligimus; sed ex parte intellectus qui est in potentia.

[31751] Iª q. 79 a. 4 ad 3
Ad tertium dicendum quod, si intellectus agens compararetur ad intellectum possibilem ut obiectum agens ad potentiam, sicut visibile in actu ad visum; sequeretur quod statim omnia intelligeremus, cum intellectus agens sit quo est omnia facere. Nunc autem non se habet ut obiectum, sed ut faciens obiecta in actu, ad quod requiritur, praeter praesentiam intellectus agentis, praesentia phantasmatum, et bona dispositio virium sensitivarum, et exercitium in huiusmodi opere; quia per unum intellectum fiunt etiam alia intellecta, sicut per terminos propositiones, et per prima principia conclusiones. Et quantum ad hoc, non differt utrum intellectus agens sit aliquid animae, vel aliquid separatum.

[31752] Iª q. 79 a. 4 ad 4
Ad quartum dicendum quod anima intellectiva est quidem actu immaterialis, sed est in potentia ad determinatas species rerum. Phantasmata autem, e converso, sunt quidem actu similitudines specierum quarundam, sed sunt potentia immaterialia. Unde nihil prohibet unam et eandem animam, inquantum est immaterialis in actu, habere aliquam virtutem per quam faciat immaterialia in actu abstrahendo a conditionibus individualis materiae, quae quidem virtus dicitur intellectus agens; et aliam virtutem receptivam huiusmodi specierum, quae dicitur intellectus possibilis, inquantum est in potentia ad huiusmodi species.

[31753] Iª q. 79 a. 4 ad 5
Ad quintum dicendum quod, cum essentia animae sit immaterialis, a supremo intellectu creata, nihil prohibet virtutem quae a supremo intellectu participatur, per quam abstrahit a materia, ab essentia ipsius procedere, sicut et alias eius potentias.


ARTICULUS 5
[31754] Iª q. 79 a. 5 arg. 1
Ad quintum sic proceditur. Videtur quod intellectus agens sit unus in omnibus. Nihil enim quod est separatum a corpore, multiplicatur secundum multiplicationem corporum. Sed intellectus agens est separatus, ut dicitur in III de anima. Ergo non multiplicatur in multis corporibus hominum, sed est unus in omnibus.

[31755] Iª q. 79 a. 5 arg. 2
Praeterea, intellectus agens facit universale, quod est unum in multis. Sed illud quod est causa unitatis, magis est unum. Ergo intellectus agens est unus in omnibus.

[31756] Iª q. 79 a. 5 arg. 3
Praeterea, omnes homines conveniunt in primis conceptionibus intellectus. His autem assentiunt per intellectum agentem. Ergo conveniunt omnes in uno intellectu agente.

[31757] Iª q. 79 a. 5 s. c.
Sed contra est quod philosophus dicit, in III de anima, quod intellectus agens est sicut lumen. Non autem est idem lumen in diversis illuminatis. Ergo non est idem intellectus agens in diversis hominibus.

[31758] Iª q. 79 a. 5 co.
Respondeo dicendum quod veritas huius quaestionis dependet ex praemissis. Si enim intellectus agens non esset aliquid animae, sed esset quaedam substantia separata, unus esset intellectus agens omnium hominum. Et hoc intelligunt qui ponunt unitatem intellectus agentis. Si autem intellectus agens sit aliquid animae, ut quaedam virtus ipsius, necesse est dicere quod sint plures intellectus agentes, secundum pluralitatem animarum, quae multiplicantur secundum multiplicationem hominum, ut supra dictum est. Non enim potest esse quod una et eadem virtus numero sit diversarum substantiarum.

[31759] Iª q. 79 a. 5 ad 1
Ad primum ergo dicendum quod philosophus probat intellectum agentem esse separatum, per hoc quod possibilis est separatus; quia, ut ipse dicit, agens est honorabilius patiente. Intellectus autem possibilis dicitur separatus, quia non est actus alicuius organi corporalis. Et secundum hunc modum etiam intellectus agens dicitur separatus, non quasi sit aliqua substantia separata.

[31760] Iª q. 79 a. 5 ad 2
Ad secundum dicendum quod intellectus agens causat universale abstrahendo a materia. Ad hoc autem non requiritur quod sit unus in omnibus habentibus intellectum, sed quod sit unus in omnibus secundum habitudinem ad omnia a quibus abstrahit universale, respectu quorum universale est unum. Et hoc competit intellectui agenti inquantum est immaterialis.

[31761] Iª q. 79 a. 5 ad 3
Ad tertium dicendum quod omnia quae sunt unius speciei, communicant in actione consequente naturam speciei, et per consequens in virtute, quae est actionis principium, non quod sit eadem numero in omnibus. Cognoscere autem prima intelligibilia est actio consequens speciem humanam. Unde oportet quod omnes homines communicent in virtute quae est principium huius actionis, et haec est virtus intellectus agentis. Non tamen oportet quod sit eadem numero in omnibus. Oportet tamen quod ab uno principio in omnibus derivetur. Et sic illa communicatio hominum in primis intelligibilibus, demonstrat unitatem intellectus separati, quem Plato comparat soli; non autem unitatem intellectus agentis, quem Aristoteles comparat lumini.


ARTICULUS 6
[31762] Iª q. 79 a. 6 arg. 1
Ad sextum sic proceditur. Videtur quod memoria non sit in parte intellectiva animae. Dicit enim Augustinus, XII de Trin., quod ad partem superiorem animae pertinent quae non sunt hominibus pecoribusque communia. Sed memoria est hominibus pecoribusque communis, dicit enim ibidem quod possunt pecora sentire per corporis sensus corporalia, et ea mandare memoriae. Ergo memoria non pertinet ad partem animae intellectivam.

[31763] Iª q. 79 a. 6 arg. 2
Praeterea, memoria praeteritorum est. Sed praeteritum dicitur secundum aliquod determinatum tempus. Memoria igitur est cognoscitiva alicuius sub determinato tempore; quod est cognoscere aliquid sub hic et nunc. Hoc autem non est intellectus, sed sensus. Memoria igitur non est in parte intellectiva, sed solum in parte sensitiva.

[31764] Iª q. 79 a. 6 arg. 3
Praeterea, in memoria conservantur species rerum quae actu non cogitantur. Sed hoc non est possibile accidere in intellectu, quia intellectus fit in actu per hoc quod informatur specie intelligibili; intellectum autem esse in actu, est ipsum intelligere in actu; et sic intellectus omnia intelligit in actu, quorum species apud se habet. Non ergo memoria est in parte intellectiva.

[31765] Iª q. 79 a. 6 s. c.
Sed contra est quod Augustinus dicit, X de Trin., quod memoria, intelligentia et voluntas sunt una mens.

[31766] Iª q. 79 a. 6 co.
Respondeo dicendum quod, cum de ratione memoriae sit conservare species rerum quae actu non apprehenduntur, hoc primum considerari oportet, utrum species intelligibiles sic in intellectu conservari possint. Posuit enim Avicenna hoc esse impossibile. In parte enim sensitiva dicebat hoc accidere, quantum ad aliquas potentias, inquantum sunt actus organorum corporalium, in quibus conservari possunt aliquae species absque actuali apprehensione. In intellectu autem, qui caret organo corporali, nihil existit nisi intelligibiliter. Unde oportet intelligi in actu illud cuius similitudo in intellectu existit. Sic ergo, secundum ipsum, quam cito aliquis actu desinit intelligere aliquam rem, desinit esse illius rei species in intellectu, sed oportet, si denuo vult illam rem intelligere, quod convertat se ad intellectum agentem, quem ponit substantiam separatam, ut ab illo effluant species intelligibiles in intellectum possibilem. Et ex exercitio et usu convertendi se ad intellectum agentem, relinquitur, secundum ipsum, quaedam habilitas in intellectu possibili convertendi se ad intellectum agentem, quam dicebat esse habitum scientiae. Secundum igitur hanc positionem, nihil conservatur in parte intellectiva, quod non actu intelligatur. Unde non poterit poni memoria in parte intellectiva, secundum hunc modum. Sed haec opinio manifeste repugnat dictis Aristotelis. Dicit enim, in III de anima, quod, cum intellectus possibilis sic fiat singula ut sciens, dicitur qui secundum actum; et quod hoc accidit cum possit operari per seipsum. Est quidem igitur et tunc potentia quodammodo; non tamen similiter ut ante addiscere aut invenire. Dicitur autem intellectus possibilis fieri singula, secundum quod recipit species singulorum. Ex hoc ergo quod recipit species intelligibilium, habet quod possit operari cum voluerit, non autem quod semper operetur, quia et tunc est quodammodo in potentia, licet aliter quam ante intelligere; eo scilicet modo quo sciens in habitu est in potentia ad considerandum in actu. Repugnat etiam praedicta positio rationi. Quod enim recipitur in aliquo, recipitur in eo secundum modum recipientis. Intellectus autem est magis stabilis naturae et immobilis, quam materia corporalis. Si ergo materia corporalis formas quas recipit, non solum tenet dum per eas agit in actu, sed etiam postquam agere per eas cessaverit; multo fortius intellectus immobiliter et inamissibiliter recipit species intelligibiles, sive a sensibilibus acceptas, sive etiam ab aliquo superiori intellectu effluxas. Sic igitur, si memoria accipiatur solum pro vi conservativa specierum, oportet dicere memoriam esse in intellectiva parte. Si vero de ratione memoriae sit quod eius obiectum sit praeteritum, ut praeteritum; memoria in parte intellectiva non erit, sed sensitiva tantum, quae est apprehensiva particularium. Praeteritum enim, ut praeteritum, cum significet esse sub determinato tempore, ad conditionem particularis pertinet.

[31767] Iª q. 79 a. 6 ad 1
Ad primum ergo dicendum quod memoria, secundum quod est conservativa specierum, non est nobis pecoribusque communis. Species enim conservantur non in parte animae sensitiva tantum, sed magis in coniuncto; cum vis memorativa sit actus organi cuiusdam. Sed intellectus secundum seipsum est conservativus specierum, praeter concomitantiam organi corporalis. Unde philosophus dicit, in III de anima, quod anima est locus specierum, non tota, sed intellectus.

[31768] Iª q. 79 a. 6 ad 2
Ad secundum dicendum quod praeteritio potest ad duo referri, scilicet ad obiectum quod cognoscitur; et ad cognitionis actum. Quae quidem duo simul coniunguntur in parte sensitiva, quae est apprehensiva alicuius per hoc quod immutatur a praesenti sensibili, unde simul animal memoratur se prius sensisse in praeterito, et se sensisse quoddam praeteritum sensibile. Sed quantum ad partem intellectivam pertinet, praeteritio accidit, et non per se convenit, ex parte obiecti intellectus. Intelligit enim intellectus hominem, inquantum est homo, homini autem, inquantum est homo, accidit vel in praesenti vel in praeterito vel in futuro esse. Ex parte vero actus, praeteritio per se accipi potest etiam in intellectu, sicut in sensu. Quia intelligere animae nostrae est quidam particularis actus, in hoc vel in illo tempore existens, secundum quod dicitur homo intelligere nunc vel heri vel cras. Et hoc non repugnat intellectualitati, quia huiusmodi intelligere, quamvis sit quoddam particulare, tamen est immaterialis actus, ut supra de intellectu dictum est; et ideo sicut intelligit seipsum intellectus, quamvis ipse sit quidam singularis intellectus, ita intelligit suum intelligere, quod est singularis actus vel in praeterito vel in praesenti vel in futuro existens. Sic igitur salvatur ratio memoriae, quantum ad hoc quod est praeteritorum, in intellectu, secundum quod intelligit se prius intellexisse, non autem secundum quod intelligit praeteritum, prout est hic et nunc.

[31769] Iª q. 79 a. 6 ad 3
Ad tertium dicendum quod species intelligibilis aliquando est in intellectu in potentia tantum, et tunc dicitur intellectus esse in potentia. Aliquando autem secundum ultimam completionem actus, et tunc intelligit actu. Aliquando medio modo se habet inter potentiam et actum, et tunc dicitur esse intellectus in habitu. Et secundum hunc modum intellectus conservat species, etiam quando actu non intelligit.


ARTICULUS 7
[31770] Iª q. 79 a. 7 arg. 1
Ad septimum sic proceditur. Videtur quod alia potentia sit memoria intellectiva, et alia intellectus. Augustinus enim, in X de Trin., ponit in mente memoriam, intelligentiam et voluntatem. Manifestum est autem quod memoria est alia potentia a voluntate. Ergo similiter est alia ab intellectu.

[31771] Iª q. 79 a. 7 arg. 2
Praeterea, eadem ratio distinctionis est potentiarum sensitivae partis et intellectivae. Sed memoria in parte sensitiva est alia potentia a sensu, ut supra dictum est. Ergo memoria partis intellectivae est alia potentia ab intellectu.

[31772] Iª q. 79 a. 7 arg. 3
Praeterea, secundum Augustinum, memoria, intelligentia et voluntas sunt sibi invicem aequalia, et unum eorum ab alio oritur. Hoc autem esse non posset, si memoria esset eadem potentia cum intellectu. Non est ergo eadem potentia.

[31773] Iª q. 79 a. 7 s. c.
Sed contra, de ratione memoriae est, quod sit thesaurus vel locus conservativus specierum. Hoc autem philosophus, in III de anima, attribuit intellectui, ut dictum est. Non ergo in parte intellectiva est alia potentia memoria ab intellectu.

[31774] Iª q. 79 a. 7 co.
Respondeo dicendum quod, sicut supra dictum est, potentiae animae distinguuntur secundum diversas rationes obiectorum; eo quod ratio cuiuslibet potentiae consistit in ordine ad id ad quod dicitur, quod est eius obiectum. Dictum est etiam supra quod, si aliqua potentia secundum propriam rationem ordinetur ad aliquod obiectum secundum communem rationem obiecti, non diversificabitur illa potentia secundum diversitates particularium differentiarum, sicut potentia visiva, quae respicit suum obiectum secundum rationem colorati, non diversificatur per diversitatem albi et nigri. Intellectus autem respicit suum obiectum secundum communem rationem entis; eo quod intellectus possibilis est quo est omnia fieri. Unde secundum nullam differentiam entium, diversificatur differentia intellectus possibilis. Diversificatur tamen potentia intellectus agentis, et intellectus possibilis, quia respectu eiusdem obiecti, aliud principium oportet esse potentiam activam, quae facit obiectum esse in actu; et aliud potentiam passivam, quae movetur ab obiecto in actu existente. Et sic potentia activa comparatur ad suum obiectum, ut ens in actu ad ens in potentia, potentia autem passiva comparatur ad suum obiectum e converso, ut ens in potentia ad ens in actu. Sic igitur nulla alia differentia potentiarum in intellectu esse potest, nisi possibilis et agentis. Unde patet quod memoria non est alia potentia ab intellectu, ad rationem enim potentiae passivae pertinet conservare, sicut et recipere.

[31775] Iª q. 79 a. 7 ad 1
Ad primum ergo dicendum quod, quamvis in III dist. I Sent. dicatur quod memoria, intelligentia et voluntas sint tres vires; tamen hoc non est secundum intentionem Augustini, qui expresse dicit in XIV de Trin., quod si accipiatur memoria, intelligentia et voluntas, secundum quod semper praesto sunt animae, sive cogitentur sive non cogitentur, ad solam memoriam pertinere videntur. Intelligentiam autem nunc dico qua intelligimus cogitantes; et eam voluntatem, sive amorem vel dilectionem, quae istam prolem parentemque coniungit. Ex quo patet quod ista tria non accipit Augustinus pro tribus potentiis; sed memoriam accipit pro habituali animae retentione, intelligentiam autem pro actu intellectus, voluntatem autem pro actu voluntatis.

[31776] Iª q. 79 a. 7 ad 2
Ad secundum dicendum quod praeteritum et praesens possunt esse propriae differentiae potentiarum sensitivarum diversificativae; non autem potentiarum intellectivarum, ratione supra dicta.

[31777] Iª q. 79 a. 7 ad 3
Ad tertium dicendum quod intelligentia oritur ex memoria, sicut actus ex habitu. Et hoc modo etiam aequatur ei; non autem sicut potentia potentiae.


ARTICULUS 8
[31778] Iª q. 79 a. 8 arg. 1
Ad octavum sic proceditur. Videtur quod ratio sit alia potentia ab intellectu. In libro enim de spiritu et anima dicitur, cum ab inferioribus ad superiora ascendere volumus, prius occurrit nobis sensus, deinde imaginatio, deinde ratio, deinde intellectus. Est ergo alia potentia ratio ab intellectu, sicut imaginatio a ratione.

[31779] Iª q. 79 a. 8 arg. 2
Praeterea, Boetius dicit, in libro de Consol., quod intellectus comparatur ad rationem sicut aeternitas ad tempus. Sed non est eiusdem virtutis esse in aeternitate et esse in tempore. Ergo non est eadem potentia ratio et intellectus.

[31780] Iª q. 79 a. 8 arg. 3
Praeterea, homo communicat cum Angelis in intellectu, cum brutis vero in sensu. Sed ratio, quae est propria hominis, qua animal rationale dicitur, est alia potentia a sensu. Ergo pari ratione est alia potentia ab intellectu, qui proprie convenit Angelis, unde et intellectuales dicuntur.

[31781] Iª q. 79 a. 8 s. c.
Sed contra est quod Augustinus dicit, III super Gen. ad Litt., quod illud quo homo irrationabilibus animalibus antecellit, est ratio, vel mens, vel intelligentia, vel si quo alio vocabulo commodius appellatur. Ratio ergo et intellectus et mens sunt una potentia.

[31782] Iª q. 79 a. 8 co.
Respondeo dicendum quod ratio et intellectus in homine non possunt esse diversae potentiae. Quod manifeste cognoscitur, si utriusque actus consideretur. Intelligere enim est simpliciter veritatem intelligibilem apprehendere. Ratiocinari autem est procedere de uno intellecto ad aliud, ad veritatem intelligibilem cognoscendam. Et ideo Angeli, qui perfecte possident, secundum modum suae naturae, cognitionem intelligibilis veritatis, non habent necesse procedere de uno ad aliud; sed simpliciter et absque discursu veritatem rerum apprehendunt, ut Dionysius dicit, VII cap. de Div. Nom. Homines autem ad intelligibilem veritatem cognoscendam perveniunt, procedendo de uno ad aliud, ut ibidem dicitur, et ideo rationales dicuntur. Patet ergo quod ratiocinari comparatur ad intelligere sicut moveri ad quiescere, vel acquirere ad habere, quorum unum est perfecti, aliud autem imperfecti. Et quia motus semper ab immobili procedit, et ad aliquid quietum terminatur; inde est quod ratiocinatio humana, secundum viam inquisitionis vel inventionis, procedit a quibusdam simpliciter intellectis, quae sunt prima principia; et rursus, in via iudicii, resolvendo redit ad prima principia, ad quae inventa examinat. Manifestum est autem quod quiescere et moveri non reducuntur ad diversas potentias, sed ad unam et eandem, etiam in naturalibus rebus, quia per eandem naturam aliquid movetur ad locum, et quiescit in loco. Multo ergo magis per eandem potentiam intelligimus et ratiocinamur. Et sic patet quod in homine eadem potentia est ratio et intellectus.

[31783] Iª q. 79 a. 8 ad 1
Ad primum ergo dicendum quod illa enumeratio fit secundum ordinem actuum, non secundum distinctionem potentiarum. Quamvis liber ille non sit magnae auctoritatis.

[31784] Iª q. 79 a. 8 ad 2
Ad secundum patet responsio ex dictis. Aeternitas enim comparatur ad tempus, sicut immobile ad mobile. Et ideo Boetius comparavit intellectum aeternitati, rationem vero tempori.

[31785] Iª q. 79 a. 8 ad 3
Ad tertium dicendum quod alia animalia sunt ita infra hominem, quod non possunt attingere ad cognoscendam veritatem, quam ratio inquirit. Homo vero attingit ad cognoscendam intelligibilem veritatem, quam Angeli cognoscunt; sed imperfecte. Et ideo vis cognoscitiva Angelorum non est alterius generis a vi cognoscitiva rationis, sed comparatur ad ipsam ut perfectum ad imperfectum.


ARTICULUS 9
[31786] Iª q. 79 a. 9 arg. 1
Ad nonum sic proceditur. Videtur quod ratio superior et inferior sint diversae potentiae. Dicit enim Augustinus, XII de Trin., quod imago Trinitatis est in superiori parte rationis, non autem in inferiori. Sed partes animae sunt ipsae eius potentiae. Ergo duae potentiae sunt ratio superior et inferior.

[31787] Iª q. 79 a. 9 arg. 2
Praeterea, nihil oritur a seipso. Sed ratio inferior oritur a superiori, et ab ea regulatur et dirigitur. Ergo ratio superior est alia potentia ab inferiori.

[31788] Iª q. 79 a. 9 arg. 3
Praeterea, philosophus dicit, in VI Ethic., quod scientificum animae quo cognoscit anima necessaria, est aliud principium et alia pars animae ab opinativo et ratiocinativo, quo cognoscit contingentia. Et hoc probat per hoc, quia ad ea quae sunt genere altera, altera genere particula animae ordinatur; contingens autem et necessarium sunt altera genere, sicut corruptibile et incorruptibile. Cum autem idem sit necessarium quod aeternum, et temporale idem quod contingens; videtur quod idem sit quod philosophus vocat scientificum, et superior pars rationis, quae secundum Augustinum intendit aeternis conspiciendis et consulendis; et quod idem sit quod philosophus vocat ratiocinativum vel opinativum, et inferior ratio, quae secundum Augustinum intendit temporalibus disponendis. Est ergo alia potentia animae ratio superior, et ratio inferior.

[31789] Iª q. 79 a. 9 arg. 4
Praeterea, Damascenus dicit quod ex imaginatione fit opinio; deinde mens, diiudicans opinionem sive vera sit sive falsa, diiudicat veritatem; unde et mens dicitur a metiendo. De quibus igitur iudicatum est iam et determinatum vere, dicitur intellectus. Sic igitur opinativum, quod est ratio inferior, est aliud a mente et intellectu, per quod potest intelligi ratio superior.

[31790] Iª q. 79 a. 9 s. c.
Sed contra est quod Augustinus dicit, XII de Trin., quod ratio superior et inferior non nisi per officia distinguuntur. Non ergo sunt duae potentiae.

[31791] Iª q. 79 a. 9 co.
Respondeo dicendum quod ratio superior et inferior, secundum quod ab Augustino accipiuntur, nullo modo duae potentiae animae esse possunt. Dicit enim quod ratio superior est quae intendit aeternis conspiciendis aut consulendis, conspiciendis quidem, secundum quod ea in seipsis speculatur; consulendis vero, secundum quod ex eis accipit regulas agendorum. Ratio vero inferior ab ipso dicitur, quae intendit temporalibus rebus. Haec autem duo, scilicet temporalia et aeterna, comparantur ad cognitionem nostram hoc modo, quod unum eorum est medium ad cognoscendum alterum. Nam secundum viam inventionis, per res temporales in cognitionem devenimus aeternorum, secundum illud apostoli, ad Rom. I, invisibilia Dei per ea quae facta sunt, intellecta, conspiciuntur, in via vero iudicii, per aeterna iam cognita de temporalibus iudicamus, et secundum rationes aeternorum temporalia disponimus. Potest autem contingere quod medium, et id ad quod per medium pervenitur, ad diversos habitus pertineant, sicut principia prima indemonstrabilia pertinent ad habitum intellectus, conclusiones vero ex his deductae ad habitum scientiae. Et ideo ex principiis geometriae contingit aliquid concludere in alia scientia, puta in perspectiva. Sed eadem potentia rationis est, ad quam pertinet et medium et ultimum. Est enim actus rationis quasi quidam motus de uno in aliud perveniens, idem autem est mobile quod pertransiens medium pertingit ad terminum. Unde una et eadem potentia rationis est ratio superior et inferior. Sed distinguuntur, secundum Augustinum, per officia actuum, et secundum diversos habitus, nam superiori rationi attribuitur sapientia, inferiori vero scientia.

[31792] Iª q. 79 a. 9 ad 1
Ad primum ergo dicendum quod secundum quamcumque rationem partitionis potest pars dici. Inquantum ergo ratio dividitur secundum diversa officia, ratio superior et inferior partitiones dicuntur, et non quia sunt diversae potentiae.

[31793] Iª q. 79 a. 9 ad 2
Ad secundum dicendum quod ratio inferior dicitur a superiori deduci, vel ab ea regulari, inquantum principia quibus utitur inferior ratio, deducuntur et diriguntur a principiis superioris rationis.

[31794] Iª q. 79 a. 9 ad 3
Ad tertium dicendum quod scientificum de quo philosophus loquitur non est idem quod ratio superior, nam necessaria scibilia inveniuntur etiam in rebus temporalibus, de quibus est scientia naturalis et mathematica. Opinativum autem et ratiocinativum in minus est quam ratio inferior, quia est contingentium tantum. Nec tamen est simpliciter dicendum quod sit alia potentia qua intellectus cognoscit necessaria, et alia qua cognoscit contingentia, quia utraque cognoscit secundum eandem rationem obiecti, scilicet secundum rationem entis et veri. Unde et necessaria, quae habent perfectum esse in veritate, perfecte cognoscit; utpote ad eorum quidditatem pertingens, per quam propria accidentia de his demonstrat. Contingentia vero imperfecte cognoscit; sicut et habent imperfectum esse et veritatem. Perfectum autem et imperfectum in actu non diversificant potentiam; sed diversificant actus quantum ad modum agendi, et per consequens principia actuum et ipsos habitus. Et ideo philosophus posuit duas particulas animae, scientificum et ratiocinativum, non quia sunt duae potentiae; sed quia distinguuntur secundum diversam aptitudinem ad recipiendum diversos habitus, quorum diversitatem ibi inquirere intendit. Contingentia enim et necessaria, etsi differant secundum propria genera, conveniunt tamen in communi ratione entis, quam respicit intellectus, ad quam diversimode se habent secundum perfectum et imperfectum.

[31795] Iª q. 79 a. 9 ad 4
Ad quartum dicendum quod illa distinctio Damasceni est secundum diversitatem actuum, non secundum diversitatem potentiarum. Opinio enim significat actum intellectus qui fertur in unam partem contradictionis cum formidine alterius. Diiudicare vero, vel mensurare, est actus intellectus applicantis principia certa ad examinationem propositorum. Et ex hoc sumitur nomen mentis. Intelligere autem est cum quadam approbatione diiudicatis inhaerere.


ARTICULUS 10
[31796] Iª q. 79 a. 10 arg. 1
Ad decimum sic proceditur. Videtur quod intelligentia sit alia potentia ab intellectu. Dicitur enim in libro de spiritu et anima, quod cum ab inferioribus ad superiora ascendere volumus, prius occurrit nobis sensus, deinde imaginatio, deinde ratio, postea intellectus, et postea intelligentia. Sed imaginatio et sensus sunt diversae potentiae. Ergo et intellectus et intelligentia.

[31797] Iª q. 79 a. 10 arg. 2
Praeterea, Boetius dicit, in V de Consol., quod ipsum hominem aliter sensus, aliter imaginatio, aliter ratio, aliter intelligentia intuetur. Sed intellectus est eadem potentia cum ratione. Ergo videtur quod intelligentia sit alia potentia quam intellectus; sicut ratio est alia potentia quam imaginatio et sensus.

[31798] Iª q. 79 a. 10 arg. 3
Praeterea, actus sunt praevii potentiis, ut dicitur in II de anima. Sed intelligentia est quidam actus ab aliis actibus qui attribuuntur intellectui divisus. Dicit enim Damascenus quod primus motus intelligentia dicitur; quae vero circa aliquid est intelligentia, intentio vocatur; quae permanens et figurans animam ad id quod intelligitur, excogitatio dicitur; excogitatio vero in eodem manens, et seipsam examinans et diiudicans, phronesis dicitur (idest sapientia); phronesis autem dilatata facit cogitationem, idest interius dispositum sermonem; ex quo aiunt provenire sermonem per linguam enarratum. Ergo videtur quod intelligentia sit quaedam specialis potentia.

[31799] Iª q. 79 a. 10 s. c.
Sed contra est quod philosophus dicit, in III de anima, quod intelligentia indivisibilium est, in quibus non est falsum. Sed huiusmodi cognoscere pertinet ad intellectum. Ergo intelligentia non est alia potentia praeter intellectum.

[31800] Iª q. 79 a. 10 co.
Respondeo dicendum quod hoc nomen intelligentia proprie significat ipsum actum intellectus qui est intelligere. In quibusdam tamen libris de Arabico translatis, substantiae separatae quas nos Angelos dicimus, intelligentiae vocantur; forte propter hoc, quod huiusmodi substantiae semper actu intelligunt. In libris tamen de Graeco translatis, dicuntur intellectus seu mentes. Sic ergo intelligentia ab intellectu non distinguitur sicut potentia a potentia; sed sicut actus a potentia. Invenitur enim talis divisio etiam a philosophis. Quandoque enim ponunt quatuor intellectus, scilicet intellectum agentem, possibilem, et in habitu, et adeptum. Quorum quatuor intellectus agens et possibilis sunt diversae potentiae; sicut et in omnibus est alia potentia activa, et alia passiva. Alia vero tria distinguuntur secundum tres status intellectus possibilis, qui quandoque est in potentia tantum, et sic dicitur possibilis; quandoque autem in actu primo, qui est scientia, et sic dicitur intellectus in habitu; quandoque autem in actu secundo, qui est considerare, et sic dicitur intellectus in actu, sive intellectus adeptus.

[31801] Iª q. 79 a. 10 ad 1
Ad primum ergo dicendum quod, si recipi debet illa auctoritas, intelligentia ponitur pro actu intellectus. Et sic dividitur contra intellectum, sicut actus contra potentiam.

[31802] Iª q. 79 a. 10 ad 2
Ad secundum dicendum quod Boetius accipit intelligentiam pro actu intellectus qui transcendit actum rationis. Unde ibidem dicit quod ratio tantum humani generis est, sicut intelligentia sola divini, proprium enim Dei est quod absque omni investigatione omnia intelligat.

[31803] Iª q. 79 a. 10 ad 3
Ad tertium dicendum quod omnes illi actus quos Damascenus enumerat, sunt unius potentiae, scilicet intellectivae. Quae primo quidem simpliciter aliquid apprehendit, et hic actus dicitur intelligentia. Secundo vero, id quod apprehendit, ordinat ad aliquid aliud cognoscendum vel operandum, et hic vocatur intentio. Dum vero persistit in inquisitione illius quod intendit, vocatur excogitatio. Dum vero id quod est excogitatum examinat ad aliqua certa, dicitur scire vel sapere; quod est phronesis, vel sapientiae, nam sapientiae est iudicare, ut dicitur in I Metaphys. Ex quo autem habet aliquid pro certo, quasi examinatum, cogitat quomodo possit illud aliis manifestare, et haec est dispositio interioris sermonis; ex qua procedit exterior locutio. Non enim omnis differentia actuum potentias diversificat; sed solum illa quae non potest reduci in idem principium, ut supra dictum est.


ARTICULUS 11
[31804] Iª q. 79 a. 11 arg. 1
Ad undecimum sic proceditur. Videtur quod intellectus speculativus et practicus sint diversae potentiae. Apprehensivum enim et motivum sunt diversa genera potentiarum, ut patet in II de anima. Sed intellectus speculativus est apprehensivus tantum, intellectus autem practicus est motivus. Ergo sunt diversae potentiae.

[31805] Iª q. 79 a. 11 arg. 2
Praeterea, diversa ratio obiecti diversificat potentiam. Sed obiectum speculativi intellectus est verum, practici autem bonum; quae differunt ratione. Ergo intellectus speculativus et practicus sunt diversae potentiae.

[31806] Iª q. 79 a. 11 arg. 3
Praeterea, in parte intellectiva intellectus practicus comparatur ad speculativum, sicut aestimativa ad imaginativam in parte sensitiva. Sed aestimativa differt ab imaginativa sicut potentia a potentia, ut supra dictum est. Ergo et intellectus practicus a speculativo.

[31807] Iª q. 79 a. 11 s. c.
Sed contra est quod dicitur in III de anima, quod intellectus speculativus per extensionem fit practicus. Una autem potentia non mutatur in aliam. Ergo intellectus speculativus et practicus non sunt diversae potentiae.

[31808] Iª q. 79 a. 11 co.
Respondeo dicendum quod intellectus practicus et speculativus non sunt diversae potentiae. Cuius ratio est quia, ut supra dictum est, id quod accidentaliter se habet ad obiecti rationem quam respicit aliqua potentia, non diversificat potentiam, accidit enim colorato quod sit homo, aut magnum aut parvum; unde omnia huiusmodi eadem visiva potentia apprehenduntur. Accidit autem alicui apprehenso per intellectum, quod ordinetur ad opus, vel non ordinetur. Secundum hoc autem differunt intellectus speculativus et practicus. Nam intellectus speculativus est, qui quod apprehendit, non ordinat ad opus, sed ad solam veritatis considerationem, practicus vero intellectus dicitur, qui hoc quod apprehendit, ordinat ad opus. Et hoc est quod philosophus dicit in III de anima, quod speculativus differt a practico, fine. Unde et a fine denominatur uterque, hic quidem speculativus, ille vero practicus, idest operativus.

[31809] Iª q. 79 a. 11 ad 1
Ad primum ergo dicendum quod intellectus practicus est motivus, non quasi exequens motum, sed quasi dirigens ad motum. Quod convenit ei secundum modum suae apprehensionis.

[31810] Iª q. 79 a. 11 ad 2
Ad secundum dicendum quod verum et bonum se invicem includunt, nam verum est quoddam bonum, alioquin non esset appetibile; et bonum est quoddam verum, alioquin non esset intelligibile. Sicut igitur obiectum appetitus potest esse verum, inquantum habet rationem boni, sicut cum aliquis appetit veritatem cognoscere; ita obiectum intellectus practici est bonum ordinabile ad opus, sub ratione veri. Intellectus enim practicus veritatem cognoscit, sicut et speculativus; sed veritatem cognitam ordinat ad opus.

[31811] Iª q. 79 a. 11 ad 3
Ad tertium dicendum quod multae differentiae diversificant sensitivas potentias, quae non diversificant potentias intellectivas, ut supra dictum est.


ARTICULUS 12
[31812] Iª q. 79 a. 12 arg. 1
Ad duodecimum sic proceditur. Videtur quod synderesis sit quaedam specialis potentia ab aliis distincta. Ea enim quae cadunt sub una divisione, videntur esse unius generis. Sed in Glossa Hieronymi Ezech. I, dividitur synderesis contra irascibilem et concupiscibilem et rationalem; quae sunt quaedam potentiae. Ergo synderesis est quaedam potentia.

[31813] Iª q. 79 a. 12 arg. 2
Praeterea, opposita sunt unius generis. Sed synderesis et sensualitas opponi videntur, quia synderesis semper inclinat ad bonum, sensualitas autem semper ad malum; unde per serpentem significatur, ut patet per Augustinum, XII de Trin. Videtur ergo quod synderesis sit potentia, sicut et sensualitas.

[31814] Iª q. 79 a. 12 arg. 3
Praeterea, Augustinus dicit, in libro de libero arbitrio, quod in naturali iudicatorio adsunt quaedam regulae et semina virtutum et vera et incommutabilia, haec autem dicimus synderesim. Cum ergo regulae incommutabiles quibus iudicamus, pertineant ad rationem secundum sui superiorem partem, ut Augustinus dicit XII de Trin.; videtur quod synderesis sit idem quod ratio. Et ita est quaedam potentia.

[31815] Iª q. 79 a. 12 s. c.
Sed contra, potentiae rationales se habent ad opposita, secundum philosophum. Synderesis autem non se habet ad opposita, sed ad bonum tantum inclinat. Ergo synderesis non est potentia. Si enim esset potentia, oporteret quod esset rationalis potentia, non enim invenitur in brutis.

[31816] Iª q. 79 a. 12 co.
Respondeo dicendum quod synderesis non est potentia, sed habitus, licet quidam posuerint synderesim esse quandam potentiam ratione altiorem; quidam vero dixerint eam esse ipsam rationem, non ut est ratio, sed ut est natura. Ad huius autem evidentiam, considerandum est quod, sicut supra dictum est, ratiocinatio hominis, cum sit quidam motus, ab intellectu progreditur aliquorum, scilicet naturaliter notorum absque investigatione rationis, sicut a quodam principio immobili, et ad intellectum etiam terminatur, inquantum iudicamus per principia per se naturaliter nota, de his quae ratiocinando invenimus. Constat autem quod, sicut ratio speculativa ratiocinatur de speculativis, ita ratio practica ratiocinatur de operabilibus. Oportet igitur naturaliter nobis esse indita, sicut principia speculabilium, ita et principia operabilium. Prima autem principia speculabilium nobis naturaliter indita, non pertinent ad aliquam specialem potentiam; sed ad quendam specialem habitum, qui dicitur intellectus principiorum, ut patet in VI Ethic. Unde et principia operabilium nobis naturaliter indita, non pertinent ad specialem potentiam; sed ad specialem habitum naturalem, quem dicimus synderesim. Unde et synderesis dicitur instigare ad bonum, et murmurare de malo, inquantum per prima principia procedimus ad inveniendum, et iudicamus inventa. Patet ergo quod synderesis non est potentia, sed habitus naturalis.

[31817] Iª q. 79 a. 12 ad 1
Ad primum ergo dicendum quod illa divisio Hieronymi attenditur secundum diversitatem actuum, non secundum diversitatem potentiarum. Diversi autem actus possunt esse unius potentiae.

[31818] Iª q. 79 a. 12 ad 2
Ad secundum dicendum quod similiter oppositio sensualitatis et synderesis attenditur secundum oppositionem actuum; non sicut diversarum specierum unius generis.

[31819] Iª q. 79 a. 12 ad 3
Ad tertium dicendum quod huiusmodi incommutabiles rationes sunt prima principia operabilium, circa quae non contingit errare; et attribuuntur rationi sicut potentiae, et synderesi sicut habitui. Unde et utroque, scilicet ratione et synderesi, naturaliter iudicamus.


ARTICULUS 13
[31820] Iª q. 79 a. 13 arg. 1
Ad tertiumdecimum sic proceditur. Videtur quod conscientia sit quaedam potentia. Dicit enim Origenes quod conscientia est spiritus corrector et paedagogus animae sociatus, quo separatur a malis et adhaeret bonis. Sed spiritus in anima nominat potentiam aliquam, vel ipsam mentem, secundum illud Ephes. IV, renovamini spiritu mentis vestrae; vel ipsam imaginationem; unde et imaginaria visio spiritualis vocatur, ut patet per Augustinum, XII super Gen. ad Litt. Est ergo conscientia quaedam potentia.

[31821] Iª q. 79 a. 13 arg. 2
Praeterea, nihil est peccati subiectum nisi potentia animae. Sed conscientia est subiectum peccati, dicitur enim ad Tit. I, de quibusdam, quod inquinatae sunt eorum mens et conscientia. Ergo videtur quod conscientia sit potentia.

[31822] Iª q. 79 a. 13 arg. 3
Praeterea, necesse est quod conscientia sit vel actus, vel habitus, vel potentia. Sed non est actus, quia non semper maneret in homine. Nec est habitus, non enim esset unum quid conscientia, sed multa; per multos enim habitus cognoscitivos dirigimur in agendis. Ergo conscientia est potentia.

[31823] Iª q. 79 a. 13 s. c.
Sed contra, conscientia deponi potest, non autem potentia. Ergo conscientia non est potentia.

[31824] Iª q. 79 a. 13 co.
Respondeo dicendum quod conscientia, proprie loquendo, non est potentia, sed actus. Et hoc patet tum ex ratione nominis, tum etiam ex his quae secundum communem usum loquendi, conscientiae attribuuntur. Conscientia enim, secundum proprietatem vocabuli, importat ordinem scientiae ad aliquid, nam conscientia dicitur cum alio scientia. Applicatio autem scientiae ad aliquid fit per aliquem actum. Unde ex ista ratione nominis patet quod conscientia sit actus. Idem autem apparet ex his quae conscientiae attribuuntur. Dicitur enim conscientia testificari, ligare vel instigare, et etiam accusare vel remordere sive reprehendere. Et haec omnia consequuntur applicationem alicuius nostrae cognitionis vel scientiae ad ea quae agimus. Quae quidem applicatio fit tripliciter. Uno modo, secundum quod recognoscimus aliquid nos fecisse vel non fecisse, secundum illud Eccle. VII, scit conscientia tua te crebro maledixisse aliis, et secundum hoc, conscientia dicitur testificari. Alio modo applicatur secundum quod per nostram conscientiam iudicamus aliquid esse faciendum vel non faciendum, et secundum hoc, dicitur conscientia instigare vel ligare. Tertio modo applicatur secundum quod per conscientiam iudicamus quod aliquid quod est factum, sit bene factum vel non bene factum, et secundum hoc, conscientia dicitur excusare vel accusare, seu remordere. Patet autem quod omnia haec consequuntur actualem applicationem scientiae ad ea quae agimus. Unde proprie loquendo, conscientia nominat actum. Quia tamen habitus est principium actus, quandoque nomen conscientiae attribuitur primo habitui naturali, scilicet synderesi, sicut Hieronymus, in Glossa Ezech. I, synderesim conscientiam nominat; et Basilius naturale iudicatorium; et Damascenus dicit quod est lex intellectus nostri. Consuetum enim est quod causae et effectus per invicem nominentur.

[31825] Iª q. 79 a. 13 ad 1
Ad primum ergo dicendum quod conscientia dicitur spiritus, secundum quod spiritus pro mente ponitur, quia est quoddam mentis dictamen.

[31826] Iª q. 79 a. 13 ad 2
Ad secundum dicendum quod inquinatio dicitur esse in conscientia, non sicut in subiecto, sed sicut cognitum in cognitione, inquantum scilicet aliquis scit se esse inquinatum.

[31827] Iª q. 79 a. 13 ad 3
Ad tertium dicendum quod actus, etsi non semper maneat in se, semper tamen manet in sua causa, quae est potentia et habitus. Habitus autem ex quibus conscientia informatur, etsi multi sint, omnes tamen efficaciam habent ab uno primo, scilicet ab habitu primorum principiorum, qui dicitur synderesis. Unde specialiter hic habitus interdum conscientia nominatur, ut supra dictum est.

